\newpage
\section{Określenie wymagań szczegółowych}		%2

\hspace{0.60cm}Rozdział niniejszy przedstawia szczegółową specyfikację funkcjonalną oraz behawioralną aplikacji mobilnej „Guess the Beat” zrealizowanej w języku Kotlin na system operacyjny Android. Określono w nim strukturę interfejsu użytkownika (UI), dokładny układ kontrolek i przycisków na poszczególnych ekranach, obsługę czujników sprzętowych, automatyczne zdarzenia systemowe, reakcje oprogramowania na sytuacje awaryjne oraz perspektywy dalszego rozwoju projektu.

\subsection{Szczegółowy opis interfejsu użytkownika i układu layoutów}

\hspace{0.60cm}Interfejs graficzny aplikacji został zaprojektowany z użyciem nowoczesnych kontenerów widoków systemu Android oraz deklaratywnego frameworka Jetpack Compose, co zapewnia pełną elastyczność i skalowalność na ekranach smartfonów o różnych proporcjach i rozdzielczościach. 

\begin{itemize}
    \item \textbf{Ekran Menu Głównego (\texttt{MainMenuScreen.kt}):} 
    Zawiera wyśrodkowane logo aplikacji, estetyczne przyciski wyboru kategorii muzycznych (np. \textit{Pop}, \textit{Rock}, \textit{Lata 80.}), przycisk uruchamiający rozgrywkę (\textit{Start Game}) oraz przycisk przejścia do ekranu ustawień dźwięku i głośności. Każdy przycisk posiada zdefiniowany stan wciśnięcia (\textit{pressed state}) oraz subtelną animację powiększenia.
    
    \item \textbf{Ekran Rozgrywki (\texttt{GameScreen.kt}):} 
    Reprezentuje główny widok gry. W jego centralnej części umieszczono wskaźnik czasu oraz duży przycisk odtwarzania próbki audio. Poniżej znajduje się pole tekstowe wyszukiwarki odpowiedzi z dynamicznymi podpowiedziami oraz umieszczony obok przycisk pomijania pojedynczego etapu (\textit{Skip}), których wysokości zostały precyzyjnie dopasowane pod kątem ergonomii obsługi. Na samym dole ekranu zaimplementowano dolny pasek nawigacyjny umożliwiający szybki powrót do menu głównego oraz przejście do panelu logowania.
    
    \item \textbf{Ekran Podsumowania (\texttt{ScoreScreen.kt}):} 
    Ekran końcowy wyświetlający podsumowanie sesji: zdobyte punkty ogółem, średni czas reakcji oraz przyciski nawigacyjne: \textit{Zagraj ponownie} oraz \textit{Powrót do menu głównego}.
\end{itemize}

\subsection{Opis zachowania aplikacji, czujników i zdarzeń automatycznych}

\hspace{0.60cm}Działanie aplikacji opiera się na ściśle zdefiniowanym cyklu życia aktywności systemu Android, asynchronicznych wątkach (\textit{Kotlin Coroutines}) oraz odczytach z wbudowanych sensorów urządzenia. Poniżej opisano kluczowe scenariusze interakcji oraz zdarzenia automatyczne:

\begin{itemize}
    \item \textbf{Zdarzenie kliknięcia „Start Game”:} Po naciśnięciu przycisku w Menu Głównym, aplikacja inicjalizuje połączenie z bazą utworów, losuje pierwszy utwór z wybranej kategorii, przygotowuje bufor pliku dźwiękowego w pamięci podręcznej i płynnie przełącza widok na ekran rozgrywki (\texttt{GameScreen}).
    \item \textbf{Stopniowe wydłużanie czasu odsłuchu i pomijanie etapu:} Rozgrywka w pojedynczej rundzie podzielona jest na etapy. Początkowo gracz słyszy bardzo krótki fragment utworu ($0,2\text{ s}$). W przypadku braku pewności co do tytułu, użytkownik może pominąć aktualny etap – zamiast odrzucać cały utwór, aplikacja wydłuża czas odsłuchu kolejnej próby (sekwencyjnie: $0,2\text{ s} \rightarrow 1\text{ s} \rightarrow 3\text{ s}$ itd.) kosztem części punktów możliwych do zdobycia.
    \item \textbf{Obsługa czujnika akcelerometru (gest odwrócenia telefonu):} Aplikacja wykorzystuje wbudowany w smartfon akcelerometr. Fizyczne (lub symulowane na emulatorze) odwrócenie urządzenia ekranem do dołu zostaje automatycznie wykryte przez system i wyzwala akcję pominięcia aktualnego etapu zgadywania, stanowiąc wygodną alternatywę dla klikania przycisku na ekranie.
    \item \textbf{Zdarzenie zatwierdzenia odpowiedzi:} Po wybraniu lub wpisaniu tekstu w pole wyszukiwarki i zatwierdzeniu, aplikacja porównuje ciąg znaków z poprawnym wzorcem. W przypadku trafienia system dolicza punkty bazowe powiększone o bonus za niewykorzystane etapy i wyświetla komunikat sukcesu.
\end{itemize}

\subsection{Zachowanie aplikacji w sytuacjach niepożądanych i awaryjnych}

\hspace{0.60cm}W celu zapewnienia wysokiej niezawodności oprogramowania, aplikacja została wyposażona w mechanizmy obsługi błędów, wyjątków oraz filtracji zakłóceń sprzętowych:

\begin{itemize}
    \item \textbf{Zabezpieczenie przed nadmierną czułością akcelerometru:} Ze względu na wysoką częstotliwość próbkowania czujnika ruchu, istnieje ryzyko wielokrotnego wyzwolenia akcji pominięcia etapu podczas jednego ruchu dłonią. Aplikacja wymaga stosowania blokady czasowej (mechanizmu \textit{debouncingu}) oraz odpowiedniej kalibracji progów osi Z, co chroni gracza przed przypadkową utratą prób.
    \item \textbf{Brak połączenia lub zasobów audio:} W sytuacji, gdy plik dźwiękowy próbki z zewnętrznego API jest uszkodzony lub niedostępny z powodu utraty połączenia sieciowego, aplikacja przechwytuje wyjątek klasy \textit{IOException}, pomija uszkodzony utwór i bezpiecznie wczytuje kolejny bez przerywania rozgrywki.
    \item \textbf{Przerwanie pracy aplikacji (np. odebranie połączenia telefonicznego):} Zgodnie z cyklem życia komponentów systemu Android, w momencie przejścia aplikacji w tło (\textit{onPause}), komponent \textit{MediaPlayer} automatycznie pauzuje odtwarzanie dźwięku, a stan gry jest zachowywany w klasie \texttt{ViewModel}.
    \item \textbf{Błędne lub puste dane wejściowe:} Próba zatwierdzenia pustego pola tekstowego jest blokowana przez interfejs, zapobiegając niepotrzebnemu odjęciu próby graczowi.
\end{itemize}

\subsection{Możliwości dalszego rozwoju oprogramowania (Perspektywy)}

\hspace{0.60cm}Architektura projektu została zaprojektowana w sposób modułowy, co umożliwia jej systematyczną rozbudowę w kolejnych wersjach oprogramowania:

\begin{itemize}
    \item \textbf{System kont użytkowników (Rejestracja i Logowanie):} Rozbudowa przycisku logowania umieszczonego na dolnym pasku nawigacyjnym o pełny moduł uwierzytelniania, umożliwiający zakładanie kont oraz zapisywanie postępów i statystyk w chmurze.
    \item \textbf{Personalizacja wyglądu i obsługa ciemnego motywu (\textit{Dark Mode}):} Dopracowanie palety barw interfejsu oraz wdrożenie pełnego wsparcia dla trybu ciemnego w Jetpack Compose, poprawiającego komfort rozgrywki w słabym oświetleniu.
    \item \textbf{Tryb wieloosobowy (Multiplayer online):} Wprowadzenie rozgrywek sieciowych w czasie rzeczywistym przy użyciu WebSockets lub Firebase, umożliwiające rywalizację graczy na odległość.
\end{itemize}